\appendixitem{Zariski 拓扑}

设 V 是一个有限维向量空间。用 \(A_{V}\) 表示 V 上的以 k 为值的多项式函数所成的代数（《代数》，第四章，§5，第 10 节，定义 4）。这是一个分次代数；它的 1 次分量是 V 的对偶 \(V^{*}\)，且 \(V^{*}\) 到 \(A_{V}\) 中的单射扩张为对称代数 \(\mathbf{S}(V^{*})\) 到 \(A_{V}\) 的一个同构（《代数》，第四章，§5，第 11 节，注记 2）。

若 \((e_{1},\ldots,e_{n})\) 是 V 的一个基，且 \((\mathrm{X}_{1},\ldots,\mathrm{X}_{n})\) 是一个未定元序列，则从 \(k[\mathrm{X}_{1},\ldots,\mathrm{X}_{n}]\) 到 \(A_{V}\) 的映射把 \(k[\mathrm{X}_{1},\ldots,\mathrm{X}_{n}]\) 的任何元素 f 变为函数

\[
\sum_ {i = 1} ^ {n} \lambda_ {i} e _ {i} \mapsto f (\lambda_ {1}, \ldots , \lambda_ {n})
\]

是代数的一个同构（《代数》，第四章，§5，第 10 节，命题 19 的推论）。

命题 1. 设 H 是从 \(A_{V}\) 到 k 的代数同态所成的集。对任何 \(x \in V\)，设 \(h_{x}\) 是同态 \(f \mapsto f(x)\)（从 \(A_{V}\) 到 k 的同态）。则映射 \(x \mapsto h_{x}\) 是从 V 到 H 的一个双射。

事实上，设 \(\mathrm{H}'\) 是从 \(k[\mathrm{X}_1,\ldots ,\mathrm{X}_n]\) 到 \(k\) 的代数同态所成的集。映射 \(\chi \mapsto (\chi (\mathrm{X}_1),\dots ,\chi (\mathrm{X}_n))\) 显然是从 \(\mathrm{H}'\) 到 \(k^n\) 的一个双射。

推论. 对任何 \(x \in V\)，设 \(\mathfrak{m}_{x} = \operatorname{Ker}(h_{x})\)。则映射 \(x \mapsto \mathfrak{m}_{x}\) 是从 V 到 \(A_{V}\) 的满足 \(A_{V}/m = k\) 的理想 m 所成的集的一个双射。

子集 \(\mathrm{F}\)（\(\mathrm{V}\) 的子集）称为闭的，如果存在一个族 \((f_i)_{i\in \mathrm{I}}\)（\(\mathrm{A_V}\) 的元素的族），使得

\[
x \in \mathrm{F} \Longleftrightarrow x \in \mathrm{V} \text {   and   } f _ {i} (x) = 0 \text {   for   all   } i \in \mathrm{I}.
\]

显然 \(\varnothing\) 与 V 是闭的，且闭集的任何交是闭的。若 F 由诸 \(f_{i}\) 的为零所定义，而 \(\mathrm{F}'\) 由诸 \(f_j'\) 的为零所定义，则 \(\mathrm{F} \cup \mathrm{F}'\) 由诸 \(f_{i}f_{j}'\) 的为零所定义，因而是闭的。于是，在 V 上存在一个拓扑，使得这个拓扑的闭集恰好是上述意义下的闭集。这个拓扑称为 V 上的 Zariski 拓扑。对任何 \(f \in \mathrm{A_V}\)，用 \(\mathrm{V}_f\) 表示那些 \(x \in \mathrm{V}\)（满足 \(f(x) \neq 0\)）所成的集；这是 V 的一个开子集。显然诸 \(\mathrm{V}_f\) 构成 Zariski 拓扑的一个基。（若 \(k\) 是一个拓扑域，则 V 的典则拓扑细于 Zariski 拓扑。）

命题 1 的推论中的映射 \(x \mapsto m_{x}\) 可以看作映射 \(\varepsilon\)（从 V 到素谱 \(\operatorname{Spec}(A_{\mathrm{V}})\)（\(A_{V}\) 的素谱）的映射）（《交换代数》，第二章，§4，第 3 节，定义 4）。直接看出，Zariski 拓扑是经 \(\varepsilon\) 取 \(\operatorname{Spec}(A_{\mathrm{V}})\) 的拓扑所得的原像。

命题 2. 赋以 Zariski 拓扑的向量空间 V 是一个不可约 Noether 空间。特别地，V 的每个非空开子集都是稠密的。

由于 \(A_{V}\) 是 Noether 的，\(\operatorname{Spec}(A_{V})\) 是 Noether 的（《交换代数》，第二章，§4，第 3 节，命题 11 的推论 7），且 Noether 空间的每个子空间都是 Noether 的（出处同前，第 2 节，命题 8）。用命题 1 的推论中的记号，诸 \(m_{x}\) 的交是 \(\{0\}\)，而 \(\{0\}\) 是 \(A_{V}\) 的一个素理想；于是 V 不可约（出处同前，第 3 节，命题 14）。

